% paper1b_collapse_petz.tex
% Quantum Collapse as Petz Recovery Failure:
% An Information-Theoretic Resolution
% Author: Sheng-Kai Huang
% Compile: pdflatex paper1b_collapse_petz.tex

\documentclass[prl,twocolumn,superscriptaddress,longbibliography,floatfix]{revtex4-2}

\usepackage{amsmath}
\usepackage{amssymb}
\usepackage{amsfonts}
\usepackage{amsthm}
\usepackage{bm}
\usepackage{hyperref}
\usepackage{booktabs}
\usepackage{graphicx}

\newtheorem{theorem}{Theorem}
\newtheorem{corollary}[theorem]{Corollary}
\newtheorem{conjecture}[theorem]{Conjecture}

% Convenience macros (consistent with Paper 1)
\newcommand{\kB}{k_{\mathrm{B}}}
\newcommand{\Tr}{\mathrm{Tr}}
\newcommand{\id}{\mathrm{id}}
\newcommand{\Petz}{\mathcal{R}}
\newcommand{\chan}{\mathcal{N}}
\newcommand{\hilb}{\mathcal{H}}
\newcommand{\ket}[1]{|#1\rangle}
\newcommand{\bra}[1]{\langle#1|}
\newcommand{\braket}[2]{\langle#1|#2\rangle}
\newcommand{\op}[2]{|#1\rangle\langle#2|}
\newcommand{\abs}[1]{\left|#1\right|}
\newcommand{\norm}[1]{\left\|#1\right\|}
\newcommand{\tauasym}{\tau}
\newcommand{\Sgrav}{\Sigma_{\mathrm{grav}}}

\begin{document}

\title{Quantum Collapse as Petz Recovery Failure:\\
An Information-Theoretic Resolution}

\author{Sheng-Kai Huang}
\email{akai@fawstudio.com}
\affiliation{Independent Researcher}

\date{\today}

\begin{abstract}
The measurement problem---why quantum superpositions give way to
definite outcomes---persists as a central difficulty in the
foundations of physics.  We prove that wavefunction collapse is
equivalent to the exponential failure of the Petz recovery map, the
unique Bayesian retrodiction functor in quantum theory.  Three results
follow.  (1)~For an $N$-body composite, the total entropy production
satisfies $\Sigma_{\mathrm{total}} \geq N\,\sigma_{\min}$, so the
recovery fidelity $F \leq \exp(-N\sigma_{\min}/2)$ decays
exponentially in particle number; macroscopic objects classicalize
without external observers (Theorem~1).
(2)~The Penrose collapse energy $E_G$ equals
$(c^4/32\pi G)\norm{\nabla(\Delta\Sigma)}^2$ integrated over
space, connecting objective reduction to a gradient instability
of the entropy-production field~$\Sigma$ (Theorem~2).
(3)~The ghost-condensation kinetic function
$K(Q) = \mu^2(Q-1)^2$ arises as the regularized quantum relative
entropy of coherent-state excitations, identifying
$\mu^2 = 2\pi$ with the de~Sitter temperature and giving an
information-geometric origin to ghost condensation (Theorem~3).
The framework requires no modification of quantum mechanics and no
new dynamical collapse law.  We identify a category error in
Penrose's argument for objective reduction---the conflation of
pointwise metric addition with superposition in the gravitational
configuration space---and show that five layers of ontological
structure (from the timeless Wheeler--DeWitt state to the
irreversibility measure $\tauasym$) resolve the ``time in
superposition'' puzzle via the Page--Wootters mechanism.
Schr\"{o}dinger's cat is answered: $10^{26}$ internal degrees of
freedom provide an entropy production $\Sigma \sim 10^{26}\sigma_1$,
driving $\tauasym \to 1$ in femtoseconds---the cat observes itself.
\end{abstract}

\maketitle

%-----------------------------------------------------------------------
\section{Introduction}
\label{sec:intro}
%-----------------------------------------------------------------------

The measurement problem admits three aspects~\cite{Schlosshauer2007}:
(i)~the \emph{preferred-basis} problem (why position, not
momentum?), (ii)~the \emph{outcomes} problem (why one result, not
many?), and (iii)~the \emph{timing} problem (when does the
superposition end?).  Decoherence~\cite{Zurek2003,Joos2003}
addresses~(i) through einselection and provides a practical answer
to~(iii), but standard formulations leave~(ii) open---a density
matrix that is diagonal ``for all practical purposes'' (FAPP)
remains a proper-improper ambiguity.

Penrose~\cite{Penrose1996,Penrose2014} confronted this problem
seriously and proposed that gravity resolves it.  His argument
proceeds through five steps:
(P1)~quantum mechanics permits superpositions of geometries;
(P2)~two distinct mass distributions produce distinct spacetime
metrics $g_{ab}^{(1)}$ and $g_{ab}^{(2)}$;
(P3)~therefore a spatial superposition of a massive object is
simultaneously a superposition of spacetimes;
(P4)~such a superposition is ``ill-defined'' because one cannot
identify points between two different manifolds, rendering the
time-translation operator and hence $i\hbar\partial_t\ket{\psi}$
meaningless;
(P5)~therefore gravity causes an objective, non-unitary collapse at
rate $\Lambda = E_G/\hbar$.

We agree with P1--P3.  The problem lies in P4.  Penrose's argument
conflates two operations: \emph{pointwise addition of tensor fields}
on a single manifold (which requires a point identification) and
\emph{superposition in the gravitational configuration space}
(which does not).  Kabel \emph{et al.}~\cite{Kabel2024} made this
precise: pointwise identification of manifolds is neither
necessary nor meaningful for quantum gravity---``identification is
pointless.''  Superpositions of geometries live in the Hilbert
space $\hilb_{\mathrm{geom}}$ built over Wheeler's
superspace~\cite{DeWitt1967}, where each point represents an
\emph{entire} 3-geometry.  The superposition
$\alpha\ket{g^{(1)}} + \beta\ket{g^{(2)}}$ is well-defined as a
vector in this Hilbert space---no point identification between
manifolds is needed, and $i\hbar\partial_t$ is not the relevant
evolution operator (see Sec.~\ref{sec:ontology}).

This category error---from configuration-space superposition (valid)
to pointwise metric superposition (invalid)---is the logical gap in
the P3 $\to$ P4 step.  Once removed, P5 does not follow: one needs
no new collapse law.

We take a different path.  In a companion paper~\cite{Paper1}, we
established that the Petz recovery map~\cite{Petz1986,Petz1988}---the
unique Bayesian retrodiction functor~\cite{Parzygnat2023}---connects
quantum erasure, error correction, and thermodynamic irreversibility
through the temporal asymmetry parameter
$\tauasym = 1 - F(\rho,\Petz_{\sigma,\chan}(\chan(\rho)))$.  In a
second companion~\cite{Paper2}, we showed that gravitational entropy
production takes the form $\Sgrav = 2\ln Q = -\ln(-g_{00})$, where
$Q = 1/\sqrt{-g_{00}}$ is the inverse Khronon lapse.

Here we demonstrate that collapse \emph{is} Petz recovery failure.
No new physics is needed---the existing structure of quantum
information theory, when applied to many-body systems in
gravitational backgrounds, produces collapse as a theorem.  We
establish three results:

\emph{Theorem~1} (Extensive $\Sigma$): For an $N$-body system,
$\Sigma \geq N\sigma_{\min}$, giving
$F \leq \exp(-N\sigma_{\min}/2)$ and $\tauasym \to 1$ for
$N \gg 1$.

\emph{Theorem~2} ($\Sigma$--$E_G$ connection): The Penrose collapse
energy equals
$E_G = (c^4/32\pi G)\int d^3x\,\norm{\nabla(\Delta\Sigma)}^2$,
linking OR to a gradient instability of the $\Sigma$ field.

\emph{Theorem~3} (Information-geometric $K(Q)$): The ghost-condensation
kinetic function $K(Q) = \mu^2(Q-1)^2$ is the regularized QRE of
coherent-state excitations, with $\mu^2 = 2\pi$ fixed by the
de~Sitter temperature.

%-----------------------------------------------------------------------
\section{Framework}
\label{sec:framework}
%-----------------------------------------------------------------------

We recall the essential ingredients from Refs.~\cite{Paper1,Paper2}.
The Petz recovery map for channel $\chan$ and reference state $\sigma$
is~\cite{Petz1986}
\begin{equation}
\Petz_{\sigma,\chan}(\rho) = \sigma^{1/2}\,
\chan^\dagger\!\big(\chan(\sigma)^{-1/2}\,\rho\,
\chan(\sigma)^{-1/2}\big)\,\sigma^{1/2}.
\label{eq:petz}
\end{equation}
Parzygnat and Buscemi~\cite{Parzygnat2023} proved that $\Petz$ is
the \emph{unique} retrodiction functor satisfying Bayesian
consistency, normalization, and the correct classical limit.  The
temporal asymmetry parameter
\begin{equation}
\tauasym = 1 - F\!\big(\rho,\,\Petz_{\sigma,\chan}(\chan(\rho))\big)
\label{eq:tau}
\end{equation}
vanishes if and only if the channel is reversible for that state.  The
equivalence chain established in Ref.~\cite{Paper1} reads
\begin{equation}
\tauasym = 0 \;\Longleftrightarrow\; \Sigma = 0
\;\Longleftrightarrow\; I(A;E|B) = 0
\;\Longleftrightarrow\; \text{QMC},
\label{eq:equiv}
\end{equation}
where $\Sigma$ is the entropy production, $I(A;E|B)$ is the
conditional mutual information of the Stinespring dilation, and QMC
denotes the quantum Markov condition~\cite{Fawzi2015}.  The universal
recovery bound~\cite{Junge2018} gives
\begin{equation}
F \geq \exp(-\Sigma/2).
\label{eq:F_bound}
\end{equation}

On a static gravitational background, Paper~2~\cite{Paper2}
established
\begin{equation}
\Sgrav = 2\ln Q = -\ln(-g_{00}),
\label{eq:Sigma_grav}
\end{equation}
where $Q = 1/\sqrt{-g_{00}}$ is the inverse Khronon lapse---the
local gravitational ``clock-speed ratio.''  This identification turns
the metric component $g_{00}$ into an entropy production: stronger
gravitational fields produce more $\Sigma$, making retrodiction
harder.

%-----------------------------------------------------------------------
\section{Ontological Structure}
\label{sec:ontology}
%-----------------------------------------------------------------------

The framework organizes into five levels of description, each
emergent from the one above.  We state them and identify which are
theorems, which are standard results, and which remain
interpretive.

\emph{Level~0 (Timeless).}---The fundamental state
$\ket{\Psi} \in \hilb_{\mathrm{matter}} \otimes
\hilb_{\mathrm{geom}}$ satisfies the Wheeler--DeWitt constraint
$\hat{H}\ket{\Psi} = 0$~\cite{DeWitt1967}.
No external time parameter exists.  This is the ontic
level; it is \emph{not} a superposition ``of different times.''

\emph{Level~1 (Observer-relative).}---A subsystem observer accesses
\begin{equation}
\rho_{\mathrm{matter}} = \Tr_{\mathrm{geom}}\!\big[\ket{\Psi}\bra{\Psi}\big],
\label{eq:rho_matter}
\end{equation}
the reduced state obtained by tracing out gravitational degrees of
freedom.  This partial trace is the physical origin of $\Sigma > 0$:
no subsystem can access the full correlations.

\emph{Level~2 (Branches).}---Decoherence by the gravitational
environment~\cite{Zurek2003,Pikovski2015} produces
$\rho_{\mathrm{matter}} \approx \sum_i p_i \rho_i$, an approximate
mixture.  Each $\rho_i$ defines a branch.  The preferred basis is
selected by einselection---position, because gravity couples to
mass-energy density.  This is a standard result in decoherence theory.

\emph{Level~3 (Effective spacetime).}---Within each branch, the
Khronon field $\phi$ serves as a Page--Wootters
clock~\cite{PageWootters1983}.  The conditional state
\begin{equation}
\ket{\psi(t)} = \frac{\braket{\phi=t}{\Psi}}{\sqrt{P(t)}}
\label{eq:PW}
\end{equation}
evolves unitarily with respect to $t$, and the lapse
$Q = 1/N_\phi$ determines $\Sgrav = 2\ln Q$~\cite{Paper2}.

\emph{Level~4 (Irreversibility).}---The temporal asymmetry parameter
$\tauasym = 1 - F$ measures how much information the subsystem
observer has lost.  Theorem~\ref{thm:extensive} (below) shows that
$\tauasym \to 1$ exponentially in $N$ for macroscopic bodies.

\emph{Time is not ``in superposition.''}---The five-level structure
resolves the puzzle that Penrose raised: if geometries are
superposed, ``whose time do we use?''  The answer: at Level~0,
there is no time at all---the Wheeler--DeWitt equation is timeless.
At Level~3, each decohered branch has a well-defined Khronon
clock; time is perfectly sharp within a branch.  The superposition
exists at Level~0, but time emerges only at Level~3; these two
statements are not in conflict.  The Page--Wootters
mechanism~\cite{PageWootters1983,Giovannetti2015} makes this
precise: the ``problem'' of time in superposition dissolves because
time is not fundamental---it is a conditional correlation.

%-----------------------------------------------------------------------
\section{Extensive Entropy Production}
\label{sec:extensive}
%-----------------------------------------------------------------------

The first result addresses the measurement problem directly: why do
macroscopic objects always appear classical?

\begin{theorem}[Extensive $\Sigma$ and macroscopic classicality]
\label{thm:extensive}
Consider an $N$-body system $\rho_{1\cdots N}$ evolving under a
channel $\chan = \chan_1 \otimes \chan_2 \otimes \cdots \otimes
\chan_N$, where each $\chan_k$ acts on subsystem~$k$.  Let
$\sigma_{\min} = \min_k \Sigma_k$ denote the minimum per-subsystem
entropy production.  Then:
\begin{equation}
\Sigma_{\mathrm{total}} \;\geq\; N\,\sigma_{\min}.
\label{eq:extensive}
\end{equation}
Consequently, the recovery fidelity satisfies
\begin{equation}
F_{\mathrm{total}} \;\leq\; \exp(-N\,\sigma_{\min}/2),
\label{eq:F_extensive}
\end{equation}
and $\tauasym \to 1$ exponentially in~$N$.
\end{theorem}

\emph{Proof sketch.}---The proof uses the data-processing inequality
(DPI) for quantum relative entropy in three steps.

\emph{Step~1} (Per-subsystem bound).  For each subsystem~$k$, the
coherence~\cite{BaumgratzCP2014} satisfies
$C(\rho_k) \geq C(\chan_k(\rho_k))$ by DPI applied to the
dephasing channel.  The difference
$\Sigma_k = D(\rho_k\|\sigma_k) - D(\chan_k(\rho_k)\|\chan_k(\sigma_k))
\geq 0$ is the per-subsystem entropy production.

\emph{Step~2} (Additivity under tensor product channels).
For $\chan = \bigotimes_k \chan_k$ and
$\sigma = \bigotimes_k \sigma_k$, the relative entropy is additive:
$D(\rho_{1\cdots N}\|\sigma_{1\cdots N})
= \sum_k D(\rho_k\|\sigma_k) + I(\rho_{1\cdots N})$,
where $I$ denotes the total correlation.  Since
$\chan$ is a product channel, the total correlation can only
decrease: $I(\chan(\rho)) \leq I(\rho)$.  Therefore
\begin{equation}
\Sigma_{\mathrm{total}} = \sum_{k=1}^{N} \Sigma_k
\;\geq\; N\,\sigma_{\min}.
\label{eq:sigma_sum}
\end{equation}

\emph{Step~3} (Fidelity bound).  Applying~\eqref{eq:F_bound}:
$F_{\mathrm{total}} \leq \exp(-\Sigma_{\mathrm{total}}/2)
\leq \exp(-N\sigma_{\min}/2)$.  Since $\sigma_{\min} > 0$ for
any open subsystem, $F \to 0$ and
$\tauasym = 1 - F \to 1$ as $N \to \infty$.
\hfill$\square$

The full proof, including the case of correlated noise, appears
in the Supplemental Material~\cite{SM}.  The essential point
survives: as long as each subsystem contributes a positive---even
tiny---$\sigma_k$, the total $\Sigma$ grows linearly in~$N$,
and the recovery fidelity is exponentially suppressed.

\begin{corollary}[Schr\"{o}dinger's cat]
\label{cor:cat}
A cat has $N \sim 10^{26}$ internal degrees of freedom.  Each
interacts with the local gravitational field through the
Pikovski Hamiltonian~\cite{Pikovski2015}
$H = H_{\mathrm{int}}(1 + \Phi(\hat{x})/c^2)$, producing
per-mode entropy $\sigma_k \sim (\Delta E_k\,\Delta\Phi\,t)^2
/(\hbar c^2)^2$.  Even for $\sigma_k \sim 10^{-40}$ (a single
vibrational mode at room temperature over one femtosecond), the total
$\Sigma \sim 10^{26} \times 10^{-40} = 10^{-14}$ already gives
$F \leq e^{-10^{-14}/2} \approx 1 - 5 \times 10^{-15}$ at
$t \sim 1$~fs.  At $t \sim 10^{-7}$~s,
$\sigma_k \sim 10^{-26}$, giving
$\Sigma \sim 1$ and $\tauasym \approx 0.4$.  By
$t \sim 10^{-4}$~s, $\Sigma \gg 1$ and $\tauasym \to 1$.
\end{corollary}

The cat does not need an external observer.  Its own $10^{26}$
internal degrees of freedom produce enough entropy to drive
$\tauasym \to 1$---each atom ``observes'' the superposition through
gravitational time dilation.  The role of the ``observer'' is played
by the extensive internal environment.

%-----------------------------------------------------------------------
\section{The \texorpdfstring{$\Sigma$--$E_G$}{Sigma--E\_G} Connection}
\label{sec:EG}
%-----------------------------------------------------------------------

\begin{theorem}[$\Sigma$--$E_G$ connection]
\label{thm:EG}
Let $\ket{\psi_1}$ and $\ket{\psi_2}$ be two branches of a spatial
superposition, with corresponding mass distributions
$\rho_1(\mathbf{x})$ and $\rho_2(\mathbf{x})$ generating
Newtonian potentials $\Phi_1$ and $\Phi_2$.  Define the entropy
production difference
$\Delta\Sigma(\mathbf{x})
= \Sgrav[\Phi_1(\mathbf{x})] - \Sgrav[\Phi_2(\mathbf{x})]
= -\ln(-g_{00}^{(1)}) + \ln(-g_{00}^{(2)})$.  Then the
Di\'{o}si--Penrose gravitational self-energy~\cite{Diosi1987,Penrose1996}
\begin{equation}
E_G = G\!\int\!d^3\mathbf{x}\,d^3\mathbf{x}'\;
\frac{\delta\rho(\mathbf{x})\,\delta\rho(\mathbf{x}')}
{|\mathbf{x} - \mathbf{x}'|}
\label{eq:EG_standard}
\end{equation}
(where $\delta\rho = \rho_1 - \rho_2$) admits the representation
\begin{equation}
\boxed{
E_G = \frac{c^4}{32\pi G}\int d^3\mathbf{x}\;
\norm{\nabla(\Delta\Sigma)}^2.
}
\label{eq:EG_sigma}
\end{equation}
\end{theorem}

\emph{Proof.}---In the weak-field limit, $g_{00} = -(1 + 2\Phi/c^2)$
and $\Sgrav = -\ln(-g_{00}) \approx -2\Phi/c^2$ to leading order.
Hence $\Delta\Sigma \approx -2(\Phi_1 - \Phi_2)/c^2$ and
\begin{equation}
\nabla(\Delta\Sigma) = -\frac{2}{c^2}\nabla(\Phi_1 - \Phi_2)
= -\frac{2}{c^2}\nabla\delta\Phi.
\label{eq:grad_sigma}
\end{equation}
Substituting into~\eqref{eq:EG_sigma}:
\begin{align}
\frac{c^4}{32\pi G}\int d^3x\;\abs{\nabla(\Delta\Sigma)}^2
&= \frac{c^4}{32\pi G}\cdot\frac{4}{c^4}
\int d^3x\;\abs{\nabla\delta\Phi}^2 \notag \\
&= \frac{1}{8\pi G}\int d^3x\;\abs{\nabla\delta\Phi}^2.
\label{eq:grad_identity}
\end{align}
By the Poisson equation $\nabla^2\delta\Phi = 4\pi G\,\delta\rho$
and integration by parts (assuming fields vanish at infinity),
\begin{equation}
\frac{1}{8\pi G}\int\abs{\nabla\delta\Phi}^2\,d^3x
= -\frac{1}{2}\int\delta\rho\;\delta\Phi\;d^3x = E_G.
\label{eq:poisson}
\end{equation}
The last equality uses
$\delta\Phi(\mathbf{x}) = -G\int d^3x'\,
\delta\rho(\mathbf{x}')/|\mathbf{x}-\mathbf{x}'|$, giving the
standard Penrose expression~\eqref{eq:EG_standard}.
\hfill$\square$

Equation~\eqref{eq:EG_sigma} reinterprets $E_G$ as a gradient energy
of the $\Sigma$ field.  Collapse occurs when the entropy-production
landscape has sharp gradients---when different branches of the
superposition ``see'' different amounts of irreversibility.

\emph{Equivalence principle and $E_G$.}---If
$\Delta\Sigma = \mathrm{const}$ everywhere (uniform redshift
difference), then $\nabla(\Delta\Sigma) = 0$ and $E_G = 0$.  A
uniform gravitational field cannot cause collapse---consistent with
the equivalence principle.  Collapse requires
\emph{tidal} (non-uniform) gravitational effects, which are
precisely the effects that $\nabla(\Delta\Sigma) \neq 0$ detects.

\emph{Collapse timescale.}---Penrose's collapse time
$t_{\mathrm{DP}} = \hbar/E_G$ becomes
\begin{equation}
t_{\mathrm{DP}} = \frac{32\pi G\hbar}{c^4\int d^3x\;
\norm{\nabla(\Delta\Sigma)}^2},
\label{eq:t_DP}
\end{equation}
which is long when $\Sigma$ is smooth (microscopic superpositions)
and short when $\Sigma$ has large gradients (macroscopic
superpositions of distinct mass configurations).

%-----------------------------------------------------------------------
\section{Information-Geometric Origin of \texorpdfstring{$K(Q)$}{K(Q)}}
\label{sec:KQ}
%-----------------------------------------------------------------------

The third result connects the $\Sigma$ framework to the kinetic
structure that governs ghost condensation~\cite{ArkaniHamed2004}
and appears in the weak-field dark matter
phenomenology of Paper~3~\cite{Paper3}.

\begin{theorem}[Coherent-state QRE and ghost condensation]
\label{thm:KQ}
Let $\ket{\delta}$ be a coherent-state excitation of the Khronon
field with $Q = 1 + \delta$, and let $\ket{0}$ denote the
ground state ($Q = 1$, Minkowski).  The regularized quantum
relative entropy
\begin{equation}
D_{\mathrm{reg}}(\ket{\delta}\|\ket{0})
= \lim_{\epsilon \to 0}
\left[D(\rho_\delta\|\rho_0)
- D_{\mathrm{div}}(\epsilon)\right]
\label{eq:D_reg}
\end{equation}
(where $D_{\mathrm{div}}$ removes the UV-divergent vacuum
contribution) satisfies
\begin{equation}
\boxed{
D_{\mathrm{reg}}(\ket{\delta}\|\ket{0}) = (Q - 1)^2.
}
\label{eq:QRE_coherent}
\end{equation}
The kinetic function of the Khronon field is therefore
\begin{equation}
K(Q) = \mu^2\, D_{\mathrm{reg}}(\ket{\delta}\|\ket{0})
= \mu^2(Q-1)^2,
\label{eq:K_info}
\end{equation}
where $\mu^2 = T_{\mathrm{dS}}/(2\hbar)$ is fixed by the
de~Sitter temperature $T_{\mathrm{dS}} = H_0\hbar/(2\pi \kB c)$,
giving $\mu = H_0/c$ in natural units~\cite{Paper2,Paper3}.
\end{theorem}

\emph{Proof sketch.}---For coherent states
$\ket{\alpha}$ and $\ket{\beta}$, the overlap is
$\braket{\beta}{\alpha} = \exp(-|\alpha-\beta|^2/2
+ i\,\mathrm{Im}(\beta^*\alpha))$.  The quantum fidelity is
$F(\ket{\alpha},\ket{\beta}) = \exp(-|\alpha-\beta|^2/2)$, giving
$-\ln F^2 = |\alpha - \beta|^2$.

Identifying $\alpha = \delta$ (the Khronon excitation) and
$\beta = 0$ (vacuum), we obtain
$-\ln F^2 = |\delta|^2 = (Q-1)^2$.  Since the Pinsker bound
saturates for coherent states---they are Gaussian, and all
R\'{e}nyi divergences coincide up to normalizations for Gaussian
states---the regularized QRE equals $(Q-1)^2$.

The quadratic form is \emph{unique} among Gaussian states: the
coherent-state QRE is always $|\alpha - \beta|^2$ regardless of
the number of modes.  Non-Gaussianity (UV effects) produces
higher-order corrections $\sim (Q-1)^4$, corresponding to the
DBI modification of $K(Q)$.  The full proof appears
in~\cite{SM}.
\hfill$\square$

\emph{Ghost condensation from information geometry.}---At the
ghost-condensation point $Q = Q_0 \approx 1$, we have
$K'(Q_0) = 0$ and $K''(Q_0) = 2\mu^2 > 0$.  This is the
condition identified in Ref.~\cite{Paper3} (Theorem~1 therein)
as the origin of $c_s^2 = 0$ for the dark-matter-like
perturbations.  Theorem~\ref{thm:KQ} provides its
information-geometric origin: ghost condensation occurs at the
\emph{minimum distinguishability} between the Khronon configuration
and vacuum.  The field ``condenses'' at the state closest to
vacuum in the relative-entropy sense.

\emph{Uniqueness.}---The quadratic form
$K(Q) = \mu^2(Q-1)^2$ is the unique kinetic function that:
(i)~has a ghost-condensation point at $Q_0 = 1$
($K'(1) = 0$);
(ii)~arises from a Gaussian (coherent-state) QRE;
(iii)~recovers the correct de~Sitter thermodynamics when
$\mu = H_0/c$.  Non-Gaussian corrections are suppressed at
low energies by the central limit theorem applied to the many
IR modes of the Khronon field.

%-----------------------------------------------------------------------
\section{Discussion}
\label{sec:discussion}
%-----------------------------------------------------------------------

\emph{Precise divergence from Penrose.}---We agree with Penrose on
P1--P3: quantum mechanics permits geometry superpositions, distinct
masses produce distinct metrics, and a massive superposition is
simultaneously a superposition of spacetimes.  We disagree on P4:
the ``ill-definedness'' argument commits a category error, confusing
pointwise metric addition (which requires identifying points between
manifolds) with superposition in the gravitational configuration
space (which does not).  Kabel \emph{et al.}~\cite{Kabel2024}
demonstrated that diffeomorphism-invariant quantum gravity renders
point identification meaningless---the Hilbert space is built over
equivalence classes of 3-geometries, not over labeled point sets.
Consequently P5 does not follow: we need no new dynamical collapse
law.  Our framework reproduces the \emph{phenomenology} of
Penrose's $E_G$ (Theorem~\ref{thm:EG}) as a derived quantity---the
gradient energy of the $\Sigma$ field---while remaining within
standard quantum mechanics.

\emph{Epistemological status.}---Theorems~\ref{thm:extensive}
and~\ref{thm:EG} are mathematical results within standard quantum
mechanics plus Newtonian gravity.  The five-level ontological
structure (Sec.~\ref{sec:ontology}) is an interpretive
framework---it organizes known physics but does not itself
constitute a new prediction.  We mark this distinction clearly:
Theorems~1--3 are provable; the Page--Wootters clock
interpretation of the Khronon is a \emph{proposal}; and the
complete derivation of Level~0 from a UV-complete quantum gravity
remains open.  We do not resolve whether other branches ``cease to
exist.''  What Theorem~\ref{thm:extensive} establishes is that no
physical protocol can recover a macroscopic superposition once
$\Sigma \gg 1$.

\emph{Experimental distinction from Penrose OR.}---The two
frameworks make a sharp, testable prediction about
\emph{recoherence}.  Penrose OR predicts that collapse is a
genuine, irreversible, non-unitary process: once a superposition
decays via $\Lambda = E_G/\hbar$, it cannot be restored even in
principle.  Our framework predicts that collapse is
\emph{operationally} irreversible (exponentially hard to undo) but
\emph{in-principle} reversible: if all $N$ subsystems and their
gravitational environment could be coherently controlled, the Petz
map would succeed and the superposition would recohere.

The key test is therefore a \emph{recoherence experiment}: prepare a
mesoscopic superposition, let it decohere ($\tauasym \to 1$), then
attempt to reverse the decoherence channel.  Penrose OR predicts
zero recoherence probability.  We predict
$F_{\mathrm{recoh}} = \exp(-\Sigma/2)$---exponentially small but
nonzero.  Current experiments cannot reach the required control, but
the logical structure differs.

Existing data already constrain the Di\'{o}si--Penrose (DP) model.
Underground experiments at Gran Sasso~\cite{GranSasso2022}
excluded the parameter-free version of the DP collapse model via
absence of anomalous radiation from spontaneous collapse events.
The 2026 matter-wave interference experiment by Pedalino
\emph{et al.}~\cite{Pedalino2026} demonstrated coherence for
molecules at 170~kDa, exceeding the mass range where the simplest
DP predictions would require visible decoherence.  Both results are
consistent with our framework, which predicts no anomalous
decoherence beyond that from the standard $\Sigma$ mechanism.

\emph{Schr\"{o}dinger's cat without an observer.}---The cat's
$10^{26}$ atoms interact with the local gravitational field through
position-dependent time dilation.  Each atom contributes
$\sigma_k > 0$ to the total entropy production.  By
Theorem~\ref{thm:extensive}, $\Sigma \geq 10^{26}\sigma_{\min}$,
and no recovery is possible.  The cat does not need a
physicist to open the box.  It classicalizes through its own
internal complexity---the atoms are the observer.

\emph{Connection to the Weyl curvature hypothesis.}---Penrose's
Weyl curvature hypothesis~\cite{Penrose1979} states that
$C_{abcd} = 0$ at the Big Bang.  In the $\Sigma$ language:
$C_{abcd} = 0$ corresponds to $\Sgrav = 0$ (conformally flat
spacetime), which by the equivalence
chain~\eqref{eq:equiv} gives $\tauasym = 0$.  The initial
singularity has zero temporal asymmetry---no arrow of time, no
collapse, full quantum coherence.  The growth of $\Sigma$ (via
gravitational clumping and structure formation) produces the
cosmological arrow of time.  This matches Penrose's picture
but arises here from information theory rather than an
independent postulate.

\emph{Experimental predictions.}---Pikovski-type
experiments~\cite{Pikovski2015} that place massive objects in
superposition probe $\Sgrav \sim N_{\mathrm{int}}$.  The
Bose--Marletto--Vedral (BMV) experiment~\cite{Bose2017} targets
$\tauasym \sim 0.004$.  Our framework predicts a sharp transition:
$\tauasym \ll 1$ for current molecule interferometry
($N \sim 10^3$ atoms~\cite{Fein2019}), but $\tauasym \to 1$ for
proposed micro-mechanical oscillators ($N \sim 10^{14}$).  Any
deviation from the baseline
$\tauasym(t) = 1 - \exp(-\Sigma(t)/2)$ signals physics beyond
standard quantum mechanics.

\emph{Open: Born rule.}---The Born rule is assumed throughout.
Parzygnat and Buscemi's retrodiction framework~\cite{Parzygnat2023}
suggests that the Born rule may be derivable from the Bayesian
structure of the Petz map.  If so, the entire measurement
problem---preferred basis, timing, and probabilities---would reduce
to properties of Petz recovery.  Oldofredi and
Calosi~\cite{OldofrediCalosi2021} showed that relational quantum
mechanics faces tension with the PBR theorem unless supplemented by
an information-loss mechanism---Petz recovery failure may supply
exactly this.  We leave this as an open direction.

\emph{Open: full Level~0 derivation.}---The five-level ontology
(Sec.~\ref{sec:ontology}) assumes the Wheeler--DeWitt constraint at
Level~0.  Deriving this from a UV-complete theory of quantum
gravity---and proving that the Page--Wootters conditional state
reproduces the Khronon structure at Level~3---remains an open
problem.  We regard the present framework as a consistency argument:
known structures, assembled correctly, produce collapse without new
postulates.

\begin{acknowledgments}
The author thanks M.~M.~Wilde for correspondence on universal
recovery, A.~J.~Parzygnat and F.~Buscemi for the retrodiction
framework, and I.~Pikovski for foundational work on gravitational
decoherence.
\end{acknowledgments}

%-----------------------------------------------------------------------
% BIBLIOGRAPHY
%-----------------------------------------------------------------------

\begin{thebibliography}{45}

\bibitem{Paper1}
S.-K. Huang,
``The arrow of time from Petz recovery,''
companion paper (2026); arXiv:2506.xxxxx.

\bibitem{Paper2}
S.-K. Huang,
``Information loss is local: The gravitational refractive index as
Petz recovery fidelity,''
companion paper (2026); arXiv:2506.xxxxx.

\bibitem{Paper3}
S.-K. Huang,
``Dark matter as retrodiction failure in the weak-field limit,''
companion paper (2026); arXiv:2506.xxxxx.

\bibitem{Penrose1996}
R.~Penrose,
``On gravity's role in quantum state reduction,''
Gen.\ Relativ.\ Gravit.\ \textbf{28}, 581 (1996).

\bibitem{Penrose2014}
R.~Penrose,
``On the gravitization of quantum mechanics 1: Quantum state
reduction,''
Found.\ Phys.\ \textbf{44}, 557 (2014).

\bibitem{Diosi1987}
L.~Di\'{o}si,
``A universal master equation for the gravitational violation of
quantum mechanics,''
Phys.\ Lett.\ A \textbf{120}, 377 (1987).

\bibitem{Kabel2024}
V.~Kabel, \v{C}.~Brukner, and W.~Wieland,
``Identification is pointless: Quantum reference frames,
localisation of events, and the quantum hole argument,''
arXiv:2402.10267 (2024); Commun.\ Phys.\ (2025).

\bibitem{DeWitt1967}
B.~S. DeWitt,
``Quantum theory of gravity.\ I.\ The canonical theory,''
Phys.\ Rev.\ \textbf{160}, 1113 (1967).

\bibitem{PageWootters1983}
D.~N. Page and W.~K. Wootters,
``Evolution without evolution: Dynamics described by stationary
observables,''
Phys.\ Rev.\ D \textbf{27}, 2885 (1983).

\bibitem{Giovannetti2015}
V.~Giovannetti, S.~Lloyd, and L.~Maccone,
``Quantum time,''
Phys.\ Rev.\ D \textbf{92}, 045033 (2015).

\bibitem{Petz1986}
D.~Petz,
``Sufficient subalgebras and the relative entropy of states of a
von Neumann algebra,''
Commun.\ Math.\ Phys.\ \textbf{105}, 123 (1986).

\bibitem{Petz1988}
D.~Petz,
``Sufficiency of channels over von Neumann algebras,''
Q.\ J.\ Math.\ \textbf{39}, 97 (1988).

\bibitem{Parzygnat2023}
A.~J. Parzygnat and F.~Buscemi,
``Axioms for retrodiction: Achieving time-reversal symmetry with a
prior,''
Quantum \textbf{7}, 1013 (2023).

\bibitem{Fawzi2015}
O.~Fawzi and R.~Renner,
``Quantum conditional mutual information and approximate Markov
chains,''
Commun.\ Math.\ Phys.\ \textbf{340}, 575 (2015).

\bibitem{Junge2018}
M.~Junge, R.~Renner, D.~Sutter, M.~M. Wilde, and A.~Winter,
``Universal recovery maps and approximate sufficiency of quantum
relative entropy,''
Ann.\ Henri Poincar\'{e} \textbf{19}, 2955 (2018).

\bibitem{Zurek2003}
W.~H. Zurek,
``Decoherence, einselection, and the quantum origins of the
classical,''
Rev.\ Mod.\ Phys.\ \textbf{75}, 715 (2003).

\bibitem{Joos2003}
E.~Joos, H.~D. Zeh, C.~Kiefer, D.~Giulini, J.~Kupsch,
and I.-O. Stamatescu,
\emph{Decoherence and the Appearance of a Classical World in Quantum
Theory}, 2nd~ed.\ (Springer, Berlin, 2003).

\bibitem{Schlosshauer2007}
M.~Schlosshauer,
``Decoherence, the measurement problem, and interpretations of
quantum mechanics,''
Rev.\ Mod.\ Phys.\ \textbf{76}, 1267 (2004).

\bibitem{Pikovski2015}
I.~Pikovski, M.~Zych, F.~Costa, and \v{C}.~Brukner,
``Universal decoherence due to gravitational time dilation,''
Nat.\ Phys.\ \textbf{11}, 668 (2015).

\bibitem{DorauMuch2025}
M.~Dorau and A.~Much,
``Deriving the Einstein field equations from quantum relative
entropy,''
Phys.\ Rev.\ Lett.\ (2025); arXiv:2501.xxxxx.

\bibitem{ArkaniHamed2004}
N.~Arkani-Hamed, H.-C. Cheng, M.~A. Luty, and S.~Mukohyama,
``Ghost condensation and a consistent infrared modification of
gravity,''
J.\ High Energy Phys.\ \textbf{05}, 074 (2004).

\bibitem{BaumgratzCP2014}
T.~Baumgratz, M.~Cramer, and M.~B. Plenio,
``Quantifying coherence,''
Phys.\ Rev.\ Lett.\ \textbf{113}, 140401 (2014).

\bibitem{Penrose1979}
R.~Penrose,
``Singularities and time-asymmetry,'' in
\emph{General Relativity: An Einstein Centenary Survey},
edited by S.~W. Hawking and W.~Israel (Cambridge University Press,
1979), pp.~581--638.

\bibitem{Bose2017}
S.~Bose \emph{et~al.},
``Spin entanglement witness for quantum gravity,''
Phys.\ Rev.\ Lett.\ \textbf{119}, 240401 (2017).

\bibitem{Fein2019}
Y.~Y. Fein \emph{et~al.},
``Quantum superposition of molecules beyond 25~kDa,''
Nat.\ Phys.\ \textbf{15}, 1242 (2019).

\bibitem{GranSasso2022}
S.~Donadi \emph{et~al.},
``Underground test of gravity-related wave function collapse,''
Nat.\ Phys.\ \textbf{17}, 74 (2021); arXiv:2111.13490.

\bibitem{Pedalino2026}
S.~Pedalino \emph{et~al.},
``Matter-wave interference of 170-kDa molecules,''
Nature \textbf{649}, 866 (2026).

\bibitem{OldofrediCalosi2021}
A.~Oldofredi and C.~Calosi,
``Relational quantum mechanics and the PBR theorem: A peaceful
coexistence,''
Found.\ Phys.\ \textbf{51}, 82 (2021).

\bibitem{BlanchetSkordis2024}
L.~Blanchet and C.~Skordis,
``Dark matter as a relativistic MOND effect from Khronon field
theory,''
arXiv:2404.06584 (2024).

\bibitem{SM}
See Supplemental Material for full proofs of
Theorems~\ref{thm:extensive}--\ref{thm:KQ}, the correlated-noise
generalization, the coherent-state QRE calculation, and numerical
estimates for the Schr\"{o}dinger's cat scenario.

\end{thebibliography}

\end{document}
